\addtocontents{toc}{\protect\Needspace{3\baselineskip}}
\originalchapter{微波背景与暴胀}\label{ch:cmb}
\originalsection{声学振荡与角谱}
复合前光子与重子强耦合. 光子提供压强, 重子增加惯性. 若 $R_b=3\rho_b/(4\rho_\gamma)$, 有效声速 $c_s=c/\sqrt{3(1+R_b)}$. 这是对混合物绝热扰动的 $\delta p/\delta\rho$ 计算, 不是重子的热声速. 引力势推动压缩, 压强推动回弹, 于是每个波数经历不同相位的声学振荡. 共动声视界 $r_s=\int_0^{t_*}c_s\dd t/a$ 与横向共动距离给角尺度 $\theta_s=r_s/D_M(z_*)$.

天空温度写为 $T(\boldsymbol n)=\bar T[1+\Theta(\boldsymbol n)]$, 展开 $\Theta=\sum a_{\ell m}Y_{\ell m}$. 统计各向同性使 $\langle a_{\ell m}a^*_{\ell'm'}\rangle=C_\ell\delta_{\ell\ell'}\delta_{mm'}$. 单一天空的估计 $\widehat C_\ell=(2\ell+1)^{-1}\sum_m|a_{\ell m}|^2$ 对 Gaussian 场的相对标准差为 $\sqrt{2/(2\ell+1)}$, 即宇宙方差. 小 $\ell$ 可用模式少, 提高仪器精度也不能消除这一统计限制.

在物质主导、大尺度、绝热扰动和忽略势随时间变化的近似下, 引力红移与本征温度共同给 $\Theta\simeq\Phi/(3c^2)$, 即普通 Sachs--Wolfe 项. 随时间变化的势还给积分项, 速度给 Doppler 项, 小尺度有扩散阻尼. 本卷解释来源与尺度, 不把完整 Einstein--Boltzmann 辐射输运求解标作已实现.

\originalsection{视界问题与加速膨胀}
在减速的辐射或物质期, 共动 Hubble 尺度 $(aH)^{-1}$ 增长. 今日遥远区域回溯到早期时, 可能没有足够时间通过信号交换达到相似温度. 暴胀的核心条件是 $\ddot a>0$, 等价于 $\dd(aH)^{-1}/\dd t<0$. 若一块原本有因果联系的区域在这一时期被拉伸, 以后可包含整个可观测宇宙. 这给视界问题一种机制, 并不自动证明某一个暴胀模型.

取 $c=\hbar=1$, 最小耦合标量场的均匀能量与压强为 $u=\dot\phi^2/2+V$、$p=\dot\phi^2/2-V$. 场方程 $\ddot\phi+3H\dot\phi+V'=0$ 与 $H^2=u/(3\bar M_{\rm Pl}^2)$ 相容, $\bar M_{\rm Pl}=(8\pi G)^{-1/2}$. 势能占优时负压强使膨胀加速. 慢滚近似忽略 $\ddot\phi$ 与动能, 得 $3H\dot\phi\simeq-V'$、$3\bar M_{\rm Pl}^2H^2\simeq V$. 其自洽条件为
\begin{equation}
\epsilon_V=\frac{\bar M_{\rm Pl}^2}2\left(\frac{V'}V\right)^2\ll1,
\qquad\left|\eta_V\right|=\left|\bar M_{\rm Pl}^2\frac{V''}V\right|\ll1.
\end{equation}
扩张倍数的对数 $N=\int H\dd t\simeq\bar M_{\rm Pl}^{-2}\int_{\phi_f}^{\phi_i}V\dd\phi/V'$; 所需 $N$ 随能标与再加热历史改变, 不能把固定「60」作为对所有模型的严格常数.

\originalsection{平直性与量子涨落的边界}
曲率偏离满足 $\Omega-1=kc^2/(a^2H^2)$. 减速期偏离放大, 加速期偏离缩小, 因而暴胀能把初始曲率压低. 暴胀结束后场能量必须转为粒子与辐射, 即再加热; 热大爆炸的初始条件并非与再加热无关.

自由轻标量在近 de Sitter 空间的量子涨落经拉伸可提供近尺度不变扰动. 在单场、规范真空、慢滚且扰动线性的假设下, 曲率扰动谱的领先阶结果为 $\mathcal P_{\mathcal R}=H^2/(8\pi^2\epsilon_H\bar M_{\rm Pl}^2)$, $\epsilon_H=-\dot H/H^2$. 这里将曲时空量子模的规范化作为高阶工具引用, 未完成规范不变量量子作用量、再加热微观过程、多场或永恒暴胀的完整论证. 这些边界不会影响前述背景动力学、热史和选定题解的自足推导.

\originalsection{声速和宇宙方差的推导}
紧耦合、绝热混合物中 $\delta\rho_b/\rho_b=3\delta\rho_\gamma/(4\rho_\gamma)$. 总压强扰动只有光子贡献 $\delta p=c^2\delta\rho_\gamma/3$, 总质量等效密度扰动 $\delta\rho=\delta\rho_\gamma[1+3\rho_b/(4\rho_\gamma)]$. 比值便是 $c_s^2=c^2/[3(1+R_b)]$. 若忽略宇宙膨胀和势缓慢变化, 对局部平面扰动的连续与 Euler 式求导消元得到 $\ddot\delta_\gamma+c_s^2k^2\delta_\gamma=\text{引力驱动项}$. 均匀解是余弦和正弦, 每个波数在最后散射时已有不同相位, 因而角谱有多峰而不是一条无结构温度曲线. 膨胀、重子加载及势变化会改变振幅和相位, 精密计算须恢复它们.

对于 Gaussian 各向同性天空, 在实球谐基中某个 $\ell$ 有 $2\ell+1$ 个独立实 Gaussian 模, 每个方差 $C_\ell$. 正规化的平方和为 $\chi^2_{2\ell+1}$, 其均值 $2\ell+1$、方差 $2(2\ell+1)$. 所以 $\operatorname{Var}(\widehat C_\ell)=2C_\ell^2/(2\ell+1)$. 这里假设完整天空、无噪声、Gaussian, 部分天空和噪声会进一步增大估计不确定性. 以复球谐直接计数时 $a_{\ell,-m}=(-1)^ma_{\ell m}^*$, 不可把 $m$ 与 $-m$ 又当作独立复变量并错误多算一倍.
